% !TEX root = ../thesis.tex
%The \introduction command is provided as a convenience.
%if you want special chapter formatting, you'll probably want to avoid using it altogether

\chapter*{Introduction}
\addcontentsline{toc}{chapter}{Introduction}
\chaptermark{Introduction}
\markboth{Introduction}{Introduction}
% The three lines above are to make sure that the headers are right, that the intro gets included in the table of contents, and that it doesn't get numbered 1 so that chapter one is 1.

The Reed College Registrar’s Office handles all logistics for scheduling at the
school: the schedule of classes, final exam scheduling, classroom scheduling,
etc. The one schedule Reed students are responsible for making is their own
class schedule each year, which proves a difficult task for some. For incoming
freshmen especially, navigating the intricacies of the class schedule in the
face of their new Reed lives can lead to some scheduling conflicts that stump
even the strongest of academic advisors. The plague of class scheduling in
general is not unique to Reed College, though there is one friction point that
is specific to the school and the freshman in particular: Humanities 110.

Humanities 110, or more notoriously HUM 110, is a graduation requirement for all
Reed students. It is a conference-style class designed to be taken by all
freshmen each year to introduce them to humanistic inquiry in preparation for
the rest of their studies at Reed. The key scheduling aspect of HUM 110 is that
it is a year-long course in a semester-based academic calendar. The Registrar
has expressed that every year they are faced with freshmen who struggle to fit
the same HUM 110 section into both semesters of their first-year schedule. When
I came to them as an aspiring CS senior thesis student with an interest in
scheduling algorithms, they inquired about a potential solution to this HUM 110
complication: is it possible to let freshmen choose all their classes for the
year except for HUM 110, and then sort them into HUM 110 sections based on the
schedules they construct? Can we put even 80\% of students into a year-long
section that doesn’t conflict with the rest of their schedules?

This thesis is presented in collaboration with the Reed College Registrar to
come up with an algorithm that assigns Reed freshmen to year-long HUM 110
sections based on the rest of their chosen class schedules. Over the course of
this thesis, we present an algorithm that stands as a hypothetical solution to
this problem. Of course, we are not the first in the field of computer science
to study scheduling algorithms, or even class scheduling more specifically.
There exists plenty of software that will assign class schedules, likely even
for more specific cases such as HUM 110. However, The Registrar informed me that
most of their scheduling is done by hand, due to the expense of these software
programs as well as the overhaul of transitioning to an entirely new
software-based system. The benefit of Reed being a relatively small school is
that the scale of their scheduling problems is small enough for the manual
strategy to be effective, if more time-consuming. We hope to create a
specialized algorithm that would be easy and free to implement and hopefully
reduce the amount of work The Registrar has in balancing freshmen into HUM 110
sections every year.

Because class scheduling is an extensively studied field, we have formulated a
very specific problem with a very specific goal to build off of that research in
contribution to the Reed community. We want to assign students to just one
year-long class amidst the rest of their chosen classes, and we want to do it in
a way that both satisfies students and maintains the philosophy behind HUM 110.
The course is meant to serve as a foundation for the lines of inquiry students
will go on to study at Reed, and as part of that it is meant to foster an
environment where students are engaging with other students who are interested
in the range of subjects Reed has to offer, from language and literature to
social sciences to hard sciences to arts. So, we develop an algorithm that
assigns students to HUM 110 sections based on their ranked preferences over the
sections, and does so in a way that balances students of different disciplines
into sections of diverse academic interests.

In this paper, we discuss the formal mathematical guarantees of our algorithm as
well as evaluate its performance against our goals and the existing class
sign-up system. Chapter 1 provides necessary background information for
understanding the mathematical analysis and construction of the algorithm, as
well as more information on the existing literature on the subject. Chapter 2
defines the formal model for our problem which allows us to prove important
features our algorithm has to accomplish our goals. Chapter 3 discusses how we
developed the algorithm and provides pseudocode outlining how it works. Chapter
4 presents the experiments we ran to test the efficacy of our algorithm as well
as what results were produced.

